\documentclass[11pt,a4paper]{article}
\usepackage[utf8]{inputenc}
\usepackage[T1]{fontenc}
\usepackage{lmodern}
\usepackage[french]{babel}
\usepackage[margin=2.0cm]{geometry}
\usepackage{microtype}
\usepackage{amsmath,amssymb}
\usepackage{booktabs}
\usepackage{array}
\usepackage[table]{xcolor}
\usepackage{enumitem}
\usepackage{eurosym}
\usepackage{graphicx}
\usepackage{caption}
\usepackage[most]{tcolorbox}
\definecolor{navy}{HTML}{1F3864}
\definecolor{bluemid}{HTML}{2E5496}
\definecolor{lightblue}{HTML}{D6E0F0}
\definecolor{accent}{HTML}{C0491F}
\captionsetup{font=small,labelfont={bf,color=navy},labelsep=period}
\graphicspath{{../figures/}}
\newtcolorbox{cle}[1][]{colback=lightblue!40,colframe=navy,boxrule=0.6pt,
  arc=2pt,left=7pt,right=7pt,top=5pt,bottom=5pt,before skip=6pt,after skip=6pt,#1}
\setlength{\parindent}{0pt}
\setlength{\parskip}{4pt}

\title{\vspace{-1.2cm}\color{navy}\bfseries Méthodologie : que faire d'un paramètre
non identifié\\[2pt]
  \large La chaîne de traitement de $W$, de la démonstration à la valeur de l'information}
\author{Kélian Kaddouri}
\date{Juillet 2026}

\begin{document}
\maketitle
\vspace{-0.8cm}

\begin{cle}
\textbf{Le problème.} Le c\oe ur du modèle est une matrice de contagion dirigée $W$ entre les
cinq piliers DORA. Or sa \emph{direction} n'est pas mesurable sur les données publiques. Le
réflexe courant serait de la poser en silence ; la méthode retenue est de la \textbf{traiter}, en
six étapes qui se répondent. Cette fiche décrit la chaîne, indépendamment du cas DORA : elle est
transposable à tout modèle dont un paramètre structurel n'est pas identifié.
\end{cle}

\section*{Première question à évacuer : la formule change-t-elle ?}

\textbf{Non. Rien de ce qui suit ne modifie une équation du modèle.} Le stress latent du pilier
$j$ de l'entité $i$ reste
\[
X_{ij} \;=\; \underbrace{\textstyle\sum_{k\neq j} W_{jk}\,X_{ik}}_{\text{contagion dirigée}}
   \;+\; \underbrace{\sqrt{\rho_{ij}}\;Y_j}_{\text{systémique}}
   \;+\; \underbrace{\sqrt{1-\rho_{ij}}\;\varepsilon_{ij}}_{\text{propre}} ,
\qquad
\mathbf{X}_i = (I-W)^{-1}(\text{syst.} + \text{propre}),
\]
la normalisation de Leontief est inchangée, l'agrégation et la mesure de risque le sont aussi.
Ce que la chaîne ci-dessous déplace n'est pas la \emph{mécanique} mais le \textbf{statut
épistémique} de $W$ : de \og{}posée par jugement\fg{} à \og{}posée, bornée, corroborée, et dont
on sait chiffrer ce qu'il manque pour la connaître\fg{}. C'est une différence de \emph{régime de
connaissance}, non de modèle, et c'est ce qui la rend sûre : aucun résultat antérieur n'est
invalidé.

\section*{Les six étapes}

\textbf{1. Séparer l'identifié du non-identifié, exactement.} Toute matrice s'écrit de façon
unique $W = S + A$, avec $S=(W+W^\top)/2$ symétrique et $A=(W-W^\top)/2$ antisymétrique, et ces
deux parts sont \emph{orthogonales} au sens de Frobenius. La donnée voit la co-occurrence de deux
piliers, donc $S$ ; elle ne voit pas qui a précédé, donc pas $A$. La séparation est nette et sans
recouvrement : l'ignorance est \textbf{localisée}, elle ne contamine pas tout le paramètre. C'est
l'étape qui rend tout le reste possible.

\textbf{2. Démontrer la non-identifiabilité, et en donner la raison.} Un constat empirique ne
suffit pas ; il faut un test valide et une explication. Le test valide sur des comptages est le
\textbf{placebo par permutation} (on détruit l'antériorité en conservant la co-occurrence, puis
on recompare l'asymétrie) : le test intuitif du \og{}temps inversé\fg{} est \emph{dégénéré} et ne
prouve rien, ce qu'il faut dire. La raison, elle, vient de la théorie des processus réversibles :
$A$ est un \textbf{courant} qui change de signe sous renversement du temps, et l'énergie de
fluctuation (la forme de Dirichlet, soit la variation de la martingale associée) ne dépend que de
$S$. Aucune statistique invariante par renversement du temps ne peut donc identifier $A$. La
limite passe ainsi du statut de constat à celui de conséquence.

\textbf{3. Borner plutôt que poser.} La seule positivité de $W$ impose
$|A_{jk}| \le S_{jk}$ : l'ensemble d'ignorance est un \textbf{pavé}, donc convexe et énumérable,
et le capital se lit en \textbf{bornes} sur cet ensemble au lieu d'être un point.

\begin{cle}[colback=accent!7,colframe=accent!70!black]
\textbf{Le piège de cette étape, à ne pas dissimuler.} Énumérer les $2^{10}$ sommets du pavé ne
borne que si les extrêmes y sont \emph{atteints}. L'argument usuel de monotonie coordonnée par
coordonnée ne s'applique \textbf{pas} : augmenter $A_{jk}$ élève $W_{jk}$ mais \emph{abaisse}
$W_{kj}$, c'est un arbitrage. On procède donc par réfutation, en tirant des points
\textbf{intérieurs} du pavé et en vérifiant qu'aucun ne sort des bornes. L'hypothèse survit à
$3\,000$ tirages, et elle est présentée comme \textbf{vérifiée numériquement}, non comme un
théorème.
\end{cle}

\textbf{4. Chercher la corroboration dans une source d'un \emph{autre genre}.} Si la
spécification requise (horodatage fin, taxonomie par domaine de contrôle, causes communes
consolidées) est introuvable dans une base \emph{agrégée}, elle peut exister \emph{à l'intérieur}
d'un autre type de document : le rapport post-incident officiel, où une enquête reconstitue la
séquence des défaillances. On code alors un corpus de tels rapports en transitions dirigées, avec
un protocole explicite, et on lui applique \textbf{le même placebo} qu'à la base agrégée. Le
compromis est inversé : peu de volume, mais une direction établie par l'enquête. C'est la qualité
du codage qui donne la puissance là où le volume ne suffisait pas.

\begin{cle}[colback=accent!7,colframe=accent!70!black]
\textbf{Le confondant de cette étape, qui ne se lève pas par le calcul.} Une enquête officielle
est \emph{mandatée} pour trouver une cause racine et des responsabilités, et conclut très souvent
à une défaillance de gouvernance. Le placebo teste la \emph{cohérence} du codage, pas son
\emph{exogénéité} : la structure trouvée peut refléter la convention d'écriture des rapports. On
le déclare, et on prévoit deux dispositifs qui, eux, peuvent le réduire : une \textbf{fiabilité
inter-juges} par second codeur en aveugle, et une \textbf{relecture par praticiens} interrogés sur
leur expérience de terrain et non sur des rapports publics.
\end{cle}

\textbf{5. Formaliser en bayésien, pour que le posterieur \emph{reproduise} l'identification.}
Par paire, on pose $\theta$ la probabilité qu'une chaîne documentée aille dans un sens, et on la
relie au coefficient par $A = S\,(2\theta-1)$, transformation qui envoie $[0,1]$ exactement sur
l'intervalle admissible $[-S,+S]$. Avec un prior Beta et les comptes du corpus pour vraisemblance,
la conjugaison est exacte : ni simulation par chaînes de Markov, ni prior d'expert. La propriété
qui fait l'intérêt du procédé est que le posterieur \textbf{reproduit de lui-même} la structure
d'identification : les paires non documentées gardent un $\theta$ uniforme, donc un coefficient
couvrant tout l'admissible, sans qu'on ait à poser cette ignorance. Deux vérifications sont
obligatoires : la \textbf{sensibilité au prior} (dont un prior \emph{sceptique}, le plus
défavorable à la thèse) et la lecture \textbf{par paire}, qui se révèle plus prudente que le test
d'ensemble et doit primer.

\textbf{6. Chiffrer ce qu'il manque, et le hiérarchiser.} Restreindre l'ensemble admissible aux
configurations compatibles avec l'information acquise mesure, en capital, ce que cette information
valait. Surtout, en fixant \emph{une seule} dépendance à la fois, on obtient un \textbf{classement}
des données manquantes par valeur d'information. La recommandation cesse d'être un v\oe u
générique : elle devient une priorité argumentée, qui ne dépend d'aucun prior.

\section*{L'atterrissage : un énoncé à trois niveaux, pas un verdict}

\begin{center}\small
\renewcommand{\arraystretch}{1.25}
\begin{tabular}{@{}p{3.0cm} p{6.9cm} p{4.6cm}@{}}
\toprule
\textbf{Objet} & \textbf{Ce que la donnée établit} & \textbf{Posture} \\
\midrule
Structure & cohérence totale du corpus, placebo franchi & \textbf{revendiquée}, confondant déclaré \\
Direction par paire & minorité de paires concluantes & \textbf{non revendiquée}, bornes conservées \\
Amplitude & rien, et une bande subsisterait même à sens tous connus & \textbf{non identifiable} ainsi \\
\bottomrule
\end{tabular}
\end{center}

\vspace{2pt}
On revendique donc un \textbf{ordre partiel} entre piliers, non une matrice. Objet plus faible,
donc bien plus difficile à contester, et \emph{suffisant} pour porter les conclusions qui
comptent : le classement des piliers et l'ordre de remédiation. Deux objets coexistent enfin dans
le rapport final, et répondent à deux questions distinctes : les \textbf{bornes}, sans prior, donc
\emph{opposables} ; le \textbf{crédible} bayésien, plus resserré, donc \emph{informatif}.

\begin{cle}
\textbf{Ce qui est transposable hors du cas DORA.} La chaîne ne doit rien à DORA : elle
s'applique dès qu'un modèle repose sur un paramètre structurel que la donnée ne fixe pas.
\emph{Localiser} l'ignorance par une décomposition orthogonale ; \emph{démontrer} la limite avec
un test valide et lui donner une raison théorique ; \emph{borner} au lieu de poser, en auditant
l'hypothèse de calcul des bornes ; \emph{corroborer} par une source d'un autre genre en déclarant
son confondant ; \emph{formaliser} de sorte que le posterieur reproduise l'identification plutôt
que de la postuler ; et \emph{chiffrer} la valeur de ce qui manque pour en faire une
recommandation hiérarchisée. La ligne rouge, tenue d'un bout à l'autre : ne jamais présenter comme
mesuré ce qui est posé, et jamais comme calibré ce qui est seulement corroboré.
\end{cle}

\smallskip
{\footnotesize\textcolor{navy}{\textbf{Sources.}} Scripts \texttt{40} (réversibilité),
\texttt{30} (bornes), \texttt{53} (corpus de post-mortems), \texttt{54} (fiabilité inter-juges),
\texttt{55} (valeur de l'information), \texttt{56} (bayésien), \texttt{57} (propriétés portables
et audit des bornes). Chapitres 9 et 10 du mémoire.}

\end{document}
